\uuid{V2El}
\titre{Matrice/Transposée/Inversibilité}
\auteur{hanani ; mzari}
\organisation{exo7}
\licence{CC BY-NC-SA 4.0}
\datecreate{2020-01-20}
\niveau{L1}
\module{Algèbre}
\chapitre{Matrice}
\sousChapitre{Inverse, méthode de Gauss}
\difficulte{2}
\theme{Inverse d'une matrice, Matrice, Transposée, Inversibilité}
\exo7qcmid{lille-2.matrices.025}

\contenu{
\begin{question}
\qclassification{108.05}
\qtags{motcle=Matrice/Transposée/Inversibilité}
Soit $ A$ une matrice inversible. On notera  $^tA$ la transposée de $A$. Quelles sont les assertions vraies ?
\begin{answers}
\good{$3A$ est inversible}.
\good{$^tA$ est inversible}.
\good{$A^tA$ est inversible}.
\bad{$A+^tA$ est inversible}.
\end{answers}
\begin{explanations} Comme $A$ est inversible, il existe une matrice $B$ telle que $AB= BA=I$, où $I$ est la matrice identité. On en déduit : $3A$ est inversible et son inverse est $\displaystyle \frac{1}{3}B$.
\vskip0mm
$^tA$ est inversible et son inverse est $ ^tB$. En effet, $I=\, ^t (AB)=\, ^tB\, ^tA$.
\vskip0mm
$A^tA$ est inversible et son inverse est $^tBB$. En effet,  $(^tBB)A^tA=\, ^tB(BA)^tA= \, ^tB^tA=\, ^t (AB)=I$.
\vskip0mm
$A+^tA$ n'est pas nécessairement inversible. Contre exemple :
$A=\left(\begin{array}{rc} 1&1\\ -1&0 \end{array}\right).$
\end{explanations}
\end{question}
}
